\documentclass[a4paper,12pt,oneside]{article}
\makeatletter
\def\input@path{{../../}}
\makeatother
\usepackage{discrete-practice}

\begin{document}

\PracticeHeader{1}{5}{Функции, отношения, частичные порядки}

\Task{Задание 1}
Дано: $y = F(x)$, где $F \subset X \times X$, $X = Y = \{1, 2, 3, 4, 5\}$.\\
Укажите:
\begin{itemize}
    \item функциональные отношения;
    \item неполностью определённые функции.
\end{itemize}

\begin{numlist}
    \item $F = \{(1,1), (2,2), (3,3), (4,4), (5,5)\}$
    \item $F = \{(1,4), (2,4), (3,4), (4,4), (4,5)\}$
    \item $F = \{(3,1), (4,5), (1,5), (2,2), (5,3)\}$
    \item $F = \{(5,1), (1,5), (2,4), (4,2)\}$
    \item $F = \{(1,1), (1,3), (3,1)\}$
    \item $F = \{(2,2), (3,3), (4,3), (5,3)\}$
\end{numlist}

\Task{Задание 2}
Покажите, что функция $h:\mathbb{Z} \to \mathbb{Z}$, заданная формулой
\[
    h(x) = x^2,
\]
не является ни инъективной, ни сюръективной.

\Task{Задание 3}
Покажите, что функция $k:\mathbb{R} \to \mathbb{R}$, заданная формулой
\[
    k(x) = 4x + 3,
\]
является и инъективной, и сюръективной.

\Task{Задание 4}
Определите, какие из функций инъективны, а какие сюръективны.
\[
\includegraphics[width=0.65\linewidth]{P5.diagram.jpg}
\]

\clearpage

\Task{Задание 5}
\subsection*{1. Укажите отношения строгого порядка:}
\begin{numlist}
    \item Иванов выше Сидорова;
    \item Лена – сестра Наташи;
    \item отрезок $a$ короче отрезка $b$;
    \item отрезок $a$ длиннее отрезка $b$ на 2 см;
    \item Васильев знает Петрова;
    \item Иванов живёт этажом выше Соколова;
    \item лыжник Ухин бежит непосредственно за Ивиным.
\end{numlist}

\subsection*{2. Укажите отношения нестрогого порядка:}
\begin{numlist}
    \item числа $a$ и $b$ не являются двузначными;
    \item точка $a$ на числовой оси находится левее точки $b$;
    \item самолёт $a$ летит не быстрее самолёта $b$;
    \item расстояние между городами равно 100 км;
    \item дом $a$ не выше дома $b$;
    \item отрезок $a$ не короче отрезка $b$;
    \item хорошее лучше плохого.
\end{numlist}

\Task{Задание 6}
Нарисуйте диаграмму Хассе для каждого из следующих частично упорядоченных множеств:
\begin{alphlist}
    \item множество $\{1,2,3,4,5,6,10,15,30\}$ с отношением «$x$ делит $y$»;
    \item множество всех подмножеств множества $\{1,2,3,4\}$ с отношением «$X$ – подмножество $Y$».
\end{alphlist}

\Task{Задание 7}
Диаграмма Хассе частичного порядка $R$ на множестве
$A = \{a,b,c,d,e,f,g,h\}$ показана на рисунке. Перечислите элементы отношения $R$.
\[
\includegraphics[width=0.60\linewidth]{P5.hasse.jpg}
\]

\end{document}
